\documentclass[10pt,comsoc]{IEEEtran}
\usepackage{cite}
\usepackage{amsmath,amssymb,amsfonts,amsthm}
\usepackage{algorithm,algorithmic}
\usepackage{graphicx}
\usepackage{textcomp}
\usepackage{xcolor}
\usepackage{booktabs}
\usepackage{multirow}
\usepackage{subfig}

\newcommand{\reals}{\mathbb{R}}
\newcommand{\norm}[1]{\left\|#1\right\|}
\newcommand{\rank}{\operatorname{rank}}

\newtheorem{theorem}{Theorem}
\newtheorem{lemma}{Lemma}
\newtheorem{definition}{Definition}
\newtheorem{proposition}{Proposition}
\newtheorem{corollary}{Corollary}
\newtheorem{remark}{Remark}

\title{Verification-Guided Dual-Path Group Sparse Representation for Few-Shot Classification}

\author{Weiye Wang, Caikou~Chen, Xiaoguang~Qiao, and Qian~Peng%
\thanks{The authors are with the College of Information and Artificial Intelligence, Yangzhou University, Yangzhou, Jiangsu 225127, China (e-mail: yzcck@yzu.edu.cn).}%
\thanks{Manuscript prepared for submission.}}

\begin{document}
\maketitle

\begin{abstract}
Sparse Representation-based Classification (SRC) degrades under small sample sizes because class-wise training sub-dictionaries become under-complete. This paper re-examines the Dual-Path Group Sparse Representation (DP-GSR) framework, which couples a forward inductive path with a reverse transductive path. Through extensive empirical analysis, we find that the reverse path---formulated as a collaborative reconstruction of training samples from the test batch---is fundamentally limited as a standalone classifier (accuracy ceiling $\approx$48\% on raw pixels), but is highly effective as a \emph{forward-prediction verifier}. We introduce a verification score that quantifies forward-reverse consistency and achieves a separation gap of 0.43 between correct and incorrect forward predictions in PCA feature spaces. Based on this score, we propose a hard-switch fusion rule with only two interpretable thresholds, and a self-training pipeline that iteratively expands the training dictionary with high-verification samples. On the Extended Yale B and AR face databases, the full pipeline improves accuracy by 3--14 percentage points over forward-only baselines in few-shot regimes, while requiring no parameter tuning when transferred across datasets. The code is available at \texttt{https://github.com/xxx/DPGSR}.
\end{abstract}

\begin{IEEEkeywords}
Sparse representation, group sparsity, transductive learning, verification, self-training, few-shot classification.
\end{IEEEkeywords}

%========================================================================
\section{Introduction}
%========================================================================

Sparse Representation-based Classification (SRC) \cite{wright2009robust} and its collaborative variants \cite{zhang2011sparse} classify a test sample by measuring how well it is reconstructed by each class-specific training sub-dictionary. This inductive paradigm works well when training dictionaries are sufficiently complete, but degrades sharply in the small-sample-size (SSS) regime: with few training samples per class, sub-dictionaries become under-complete, residual-based decisions become ill-conditioned, and test samples are often reconstructed by irrelevant atoms.

To address this, Dual-Path Group Sparse Representation (DP-GSR) \cite{wang2024dual} introduces a reverse transductive path: the test batch is used as a dictionary to reconstruct the labeled training set, producing a reverse coefficient matrix whose row-wise sparsity pattern provides class-aware transductive evidence. The original DP-GSR fuses forward residual evidence and reverse activation evidence through a continuous mixture weight $\alpha_j = 1 - S_G(\mathbf{g}_j)$, where $S_G$ is the Hoyer sparsity index computed on the forward class-energy vector.

In this paper, we conduct the first systematic empirical investigation of the DP-GSR framework across feature spaces, sample sizes, and lighting conditions. Our central finding is that the reverse path is \textbf{not} effective as an independent classifier---its standalone accuracy ceiling is approximately 48\% on raw pixels and 74\% on PCA features---but is highly effective as a \textbf{forward-prediction verifier}. Specifically:

\begin{itemize}
\item We define a \textit{verification score} $v_j$ measuring the fraction of reverse coefficient energy concentrated on the forward-predicted class. On PCA-50 features, $v_j$ averages 0.79 when forward is correct versus 0.36 when forward is wrong---a separation gap of 0.43, which is 20--40 times larger than any single-path sparsity or residual signal.
\item Based on $v_j$, we replace the original continuous $\alpha_j$ mixture with a \textit{hard-switch fusion} rule ($\alpha_j \in \{0,1\}$) governed by only two thresholds. This eliminates the scale-mismatch problem inherent in convex combinations of heterogeneous evidence scores.
\item We develop a \textit{self-training pipeline} that uses high-verification samples to iteratively expand the training dictionary. The pipeline is parameterized by a single integer $R_{\max}$ (number of self-training rounds), with $R_{\max}=1$ recovering single-pass fusion and $R_{\max}=\infty$ running to convergence.
\end{itemize}

Experiments on Extended Yale B and AR face databases demonstrate consistent improvements. On Yale B under PCA-50 features with 38 subjects and 8 training samples per class, the full pipeline achieves 80.6\% accuracy versus 74.3\% forward-only (gain of +6.3\%). On AR with only 2 training samples per class, self-training improves accuracy from 68.7\% to 82.4\% (gain of +13.8\%). The method transfers across datasets without parameter retuning. On cross-lighting evaluation (training on one illumination subset, testing on all others), all 25 tested configurations show positive fusion gains, with the largest exceeding +26\%.

%========================================================================
\section{Related Work}
%========================================================================

\subsection{Sparse and Collaborative Representation}

SRC \cite{wright2009robust} classifies a query by identifying the class that yields the minimum $\ell_1$-regularized reconstruction residual. CRC \cite{zhang2011sparse} replaces $\ell_1$ with $\ell_2$ regularization, trading sparsity for efficiency while retaining competitive accuracy. Structured extensions such as Group Sparse CRC \cite{elhamifar2013robust} and Probabilistic CRC \cite{cai2016probabilistic} improve discriminability through block-sparse or probabilistic formulations. However, all these methods are purely inductive: they process each test sample independently against a fixed training dictionary and ignore the distributional information available in the test batch.

\subsection{Transductive and Self-Training Methods}

Transductive learning leverages unlabeled test data during inference. Low-Rank Representation \cite{liu2013low} and Sparse Subspace Clustering \cite{elhamifar2013ssc} capture global data structure through self-expression, while graph-based label propagation \cite{zhu2003semi, zhou2004learning} diffuses labels through a similarity graph. In few-shot learning, transductive prototype refinement \cite{liu2018learning} and information maximization \cite{boudiaf2020information} use query batches to adjust class prototypes or sharpen decision boundaries. Self-training \cite{scudder1965probability} iteratively adds high-confidence pseudo-labeled samples to the training set; its effectiveness depends critically on the quality of the confidence estimator. Our contribution uses the reverse-path verification score as a principled, training-free confidence estimator for self-training.

\subsection{Positioning of This Work}

DP-GSR differs from the above in that it preserves the SRC class-wise residual decision while introducing a reverse reconstruction check. Unlike global structure methods (LRR, SSC), it produces per-sample class evidence. Unlike graph methods, it operates on reconstruction energy rather than similarity graphs. Unlike deep transductive few-shot methods, it requires no pretraining, meta-learning, or complex optimization. Our key departure from the original DP-GSR \cite{wang2024dual} is replacing the continuous sparsity-guided $\alpha_j$ with a verification-gated hard switch, and introducing the self-training extension.

%========================================================================
\section{Forward and Reverse Group-Sparse Coding}
%========================================================================

\subsection{Notation and Problem Setup}

Let $\mathbf{A} \in \reals^{d \times N}$ be the training set of $N$ samples from $K$ classes, with $\Omega_k \subset \{1,\dots,N\}$ indexing the $n_k = |\Omega_k|$ samples of class $k$. Let $\mathbf{Y} \in \reals^{d \times M}$ be a batch of $M$ unlabeled test samples. The goal is to predict the class label of each $\mathbf{y}_{\cdot,j}$.

\subsection{Forward Path: Inductive Group-Sparse Coding}

The forward path reconstructs each test sample as a sparse linear combination of training atoms, with a column-wise group-sparsity penalty encouraging each query to be represented by few classes:

\begin{equation}
\label{eq:fw_opt}
\mathbf{X}^\star = \arg\min_{\mathbf{X}} \left\{ \|\mathbf{Y} - \mathbf{AX}\|_F^2 + \lambda_1 \|\mathbf{X}\|_G \right\},
\end{equation}
where $\|\mathbf{X}\|_G = \sum_{j=1}^M \sum_{k=1}^K \|\mathbf{x}_{k,j}\|_2$, and $\mathbf{x}_{k,j} \in \reals^{n_k}$ is the coefficient block of test sample $j$ on class $k$.

The problem is solved via ADMM. Introducing $\mathbf{X}=\mathbf{Q}$ with multiplier $\mathbf{W}$ and penalty $\eta$:

\begin{equation}
\begin{aligned}
\mathbf{X}^{(t+1)} &= (2\mathbf{A}^\top\mathbf{A} + \eta\mathbf{I})^{-1}(2\mathbf{A}^\top\mathbf{Y} + \eta\mathbf{Q}^{(t)} - \mathbf{W}^{(t)}),\\
\mathbf{q}_{k,j}^{(t+1)} &= \max\!\big(0, 1 - \tfrac{\lambda_1}{\eta\|\mathbf{u}_{k,j}\|_2}\big)\,\mathbf{u}_{k,j},\\
\mathbf{W}^{(t+1)} &= \mathbf{W}^{(t)} + \eta(\mathbf{X}^{(t+1)} - \mathbf{Q}^{(t+1)}),
\end{aligned}
\end{equation}
with $\eta \leftarrow \min(1.2\eta, 10^{20})$, $\eta^0 = 1.1$, stopping when $\|\mathbf{X}-\mathbf{Q}\|_\infty < 10^{-4}$.

The forward posterior proxy is obtained via temperature-controlled softmax over class residuals $r_k(\mathbf{y}_{\cdot,j}) = \|\mathbf{y}_{\cdot,j} - \mathbf{A}_k\mathbf{x}_{k,j}^\star\|_2^2$:

\begin{equation}
P_F(c_k|\mathbf{y}_{\cdot,j}) = \frac{\exp(-r_k/\tau_F)}{\sum_i \exp(-r_i/\tau_F)},\quad \tau_F = \text{mean}(r_1,\dots,r_K).
\end{equation}

\subsection{Reverse Path: Transductive Group-Sparse Coding}

The reverse path uses the test batch $\mathbf{Y}$ as a dictionary to reconstruct the training set $\mathbf{A}$:

\begin{equation}
\label{eq:rv_opt}
\mathbf{Z}^\star = \arg\min_{\mathbf{Z}} \left\{ \|\mathbf{A} - \mathbf{YZ}\|_F^2 + \lambda_2 \|\mathbf{Z}\|_R \right\},
\end{equation}
where $\|\mathbf{Z}\|_R = \sum_{j=1}^M \sum_{k=1}^K \|\mathbf{z}_{j,\Omega_k}\|_2$, with $\mathbf{z}_{j,\Omega_k} \in \reals^{1 \times n_k}$ denoting test sample $j$'s coefficients on class $k$'s training samples.

ADMM solution is symmetric to the forward path, with row-wise group soft-thresholding. The reverse activation strength is:

\begin{equation}
a_k(\mathbf{y}_{\cdot,j}) = \|\mathbf{z}_{j,\Omega_k}^\star\|_2,
\end{equation}

and the class-balanced normalized activation is:

\begin{equation}
\tilde{s}_k(\mathbf{y}_{\cdot,j}) = \frac{a_k(\mathbf{y}_{\cdot,j})}{\sqrt{n_k} \cdot \|\mathbf{A}_k\|_F}.
\end{equation}

The reverse posterior proxy is obtained by normalization: $P_R(c_k|\mathbf{y}_{\cdot,j}) = \tilde{s}_k / \sum_i \tilde{s}_i$.

\textbf{Support for variable class sizes.} Both \texttt{Forward.m} and \texttt{Reverse.m} accept either a scalar $n_t$ (equal class sizes) or a label vector $\mathbf{l}_A$. When a label vector is provided, group assignments use explicit label indexing rather than assuming fixed-size blocks, enabling the self-training pipeline where pseudo-labeled samples produce unbalanced classes.

%========================================================================
\section{Verification-Gated Fusion and Self-Training}
%========================================================================

\subsection{The Limits of Reverse Classification}

The reverse optimization solves a \emph{collaborative} reconstruction: all $M$ test atoms jointly reconstruct $N$ training atoms. Each row $\mathbf{z}_{j,\cdot}^\star$ reflects the $j$-th atom's contribution share in this global accounting, not an independent per-sample discriminative representation. Table~\ref{tab:rev_limits} shows that $\arg\max_k \tilde{s}_k$ consistently underperforms the forward path, and alternative aggregation schemes (top-$K$ training-sample voting, $\ell_1$ norm, sparse post-processing) provide no improvement.

\begin{table}[t]
\centering
\caption{Reverse path accuracy ceiling. ``Voting'' = top-10 training samples by $|z_{ji}|$.}
\label{tab:rev_limits}
\begin{tabular}{lccc}
\hline
\textbf{Feature} & \textbf{Forward} & \textbf{Reverse (argmax)} & \textbf{Reverse (Voting)} \\ \hline
Yale B pixels & 88.6\% & 47.3\% & 30.6\% \\
Yale B PCA-50 & 74.3\% & 74.1\% & -- \\
Yale B Gabor & 79.9\% & 46.4\% & -- \\
AR PCA-100 & 71--94\% & 56--89\% & -- \\ \hline
\end{tabular}
\end{table}

\subsection{Verification Score}

While the reverse path fails as a classifier, it excels at answering a simpler question: \emph{does the reverse representation support the forward prediction?} Let $\hat{k}_F = \arg\max_k P_F(c_k|\mathbf{y}_{\cdot,j})$. We define:

\begin{equation}
\label{eq:verif}
v_j = \frac{\|\mathbf{z}_{j, \Omega_{\hat{k}_F}}^\star\|_2}{\|\mathbf{z}_{j,\cdot}^\star\|_2}.
\end{equation}

$v_j \in [0,1]$ is the fraction of reverse energy concentrated on the forward-predicted class. Table~\ref{tab:verif_gap} reports its discriminative power.

\begin{table}[t]
\centering
\caption{Verification score distribution (Yale B, $K=38$, $n_k=8$).}
\label{tab:verif_gap}
\begin{tabular}{lccc}
\hline
\textbf{Feature} & \textbf{Forward correct} & \textbf{Forward wrong} & \textbf{Gap} \\ \hline
Raw pixels & 0.49 & 0.30 & 0.19 \\
PCA-50 & \textbf{0.79} & \textbf{0.36} & \textbf{0.43} \\ \hline
\end{tabular}
\end{table}

In PCA-50, $v_j$ achieves a separation gap of 0.43---an order of magnitude larger than forward sparsity concentration (gap $\approx$ 0.01) or residual contrast (gap $\approx$ 0.01). In raw pixel space, the gap shrinks to 0.19, consistent with the observation that fusion gains are minimal there.

\subsection{Hard-Switch Fusion}

Let $\Delta P_R^{(j)} = P_R^{(1)} - P_R^{(2)}$ be the reverse posterior margin. The fusion rule is:

\begin{equation}
\label{eq:hard_fusion}
\alpha_j = \begin{cases}
0, & v_j > 0.65,\\
1, & v_j \leq 0.65 \;\land\; \Delta P_R^{(j)} > 0.05,\\
0, & \text{otherwise}.
\end{cases}
\end{equation}

The final prediction is $\hat{l}_j = \arg\max_k [(1-\alpha_j)\mathbf{p}_F^{(j)} + \alpha_j\mathbf{p}_R^{(j)}]$. Since $\alpha_j \in \{0,1\}$, each sample is classified either by forward alone (high verification, or both unreliable) or by reverse alone (low verification + reverse confident).

This hard-switch design has several advantages over the original continuous $\alpha_j = 1 - S_G(\mathbf{g}_j)$: (i) it avoids scale mismatch between heterogeneous evidence scores, (ii) it uses information from \emph{both} paths (forward for the prediction, reverse for the verification), and (iii) it requires only two interpretable thresholds instead of a sparsity index with unclear calibration.

\begin{proposition}
When $\theta \to 0$ or $\delta \to 1$, the rule degenerates to forward-only SRC. When $\theta \to 1$ and $\delta \to 0$, it degenerates to reverse-only.
\end{proposition}

\subsection{Self-Training with Verification Guidance}

Samples with $v_j > 0.65$ are those where forward and reverse strongly agree. As shown in Section~V, their forward-predicted labels are $>95\%$ accurate on Yale B PCA-50. We therefore add them to the training set with pseudo-labels:

\begin{algorithm}[t]
\caption{Verification-Guided Self-Training DP-GSR}
\label{alg:main}
\begin{algorithmic}[1]
\REQUIRE $\mathbf{A}$, $\mathbf{l}_A$, $\mathbf{Y}$, $\lambda_1$, $\lambda_2$, $\theta=0.65$, $\delta=0.05$, $R_{\max}$.
\ENSURE $\hat{\mathbf{l}} \in \{1,\dots,K\}^M$.
\STATE $\mathcal{U} \leftarrow \{1,\dots,M\}$, $\hat{\mathbf{l}} \leftarrow \mathbf{0}$
\FOR{$r = 1$ to $R_{\max}$}
    \STATE Solve forward (\ref{eq:fw_opt}) and reverse (\ref{eq:rv_opt}) on $(\mathbf{A}, \mathbf{Y}_{:,\mathcal{U}})$.
    \STATE Compute $v_j$ $\forall j\in\mathcal{U}$ via (\ref{eq:verif}).
    \STATE $\mathcal{C} \leftarrow \{j \in \mathcal{U} \mid v_j > \theta\}$
    \IF{$|\mathcal{C}| = 0$} \STATE \textbf{break} \ENDIF
    \FOR{$j \in \mathcal{C}$}
        \STATE $\hat{k}_F \leftarrow \arg\max_k P_F(c_k|\mathbf{y}_{\cdot,j})$
        \STATE $\mathbf{A} \leftarrow [\mathbf{A}, \mathbf{y}_{\cdot,j}]$, $\mathbf{l}_A \leftarrow [\mathbf{l}_A, \hat{k}_F]$
        \STATE $\hat{l}_j \leftarrow \hat{k}_F$, $\mathcal{U} \leftarrow \mathcal{U} \setminus \{j\}$
    \ENDFOR
\ENDFOR
\STATE // Final pass: reverse on full $\mathbf{Y}$, hard-switch fusion on $\mathcal{U}$
\STATE Solve reverse on $(\mathbf{A}, \mathbf{Y})$; apply (\ref{eq:hard_fusion}) to $j\in\mathcal{U}$.
\RETURN $\hat{\mathbf{l}}$
\end{algorithmic}
\end{algorithm}

The self-training loop expands the training dictionary with verified samples. The final pass uses the full test batch as the reverse dictionary---critical because the reverse path requires $M \geq N$ for stable reconstruction. $R_{\max}=1$ recovers single-pass verification-gated fusion; $R_{\max}=\infty$ runs to convergence (typically 3--7 rounds).

%========================================================================
\section{Experiments}
%========================================================================

\subsection{Setup}

\textbf{Datasets.} Extended Yale B \cite{georghiades2001from} (38 subjects, 64 images each, 5 lighting subsets) and AR Face \cite{martinez1998ar} (100 subjects, 26 images each). Images are projected to $d=50$ (Yale B) or $d=100$ (AR) via PCA, then $\ell_2$-normalized.

\textbf{Parameters.} All experiments use $\lambda_1=10^{-1}$, $\lambda_2=10^{-2}$, $\theta=0.65$, $\delta=0.05$. These were tuned once on Yale B PCA-50 and transferred directly to AR without modification. ADMM: $\eta^0=\rho^0=1.1$, growth factor $1.2$, $\epsilon=10^{-4}$, max 500 iterations.

\textbf{Metrics.} We report forward-only accuracy (Fwd), hard-switch fusion accuracy with $R_{\max}=1$ (Fus), and self-training accuracy with $R_{\max}=\infty$ (ST). Gain is relative to forward-only.

\subsection{Main Results: Random Splits}

\begin{table}[t]
\centering
\caption{Results on Yale B (PCA-50, $K=38$, $n_k=8$, random splits, 5 trials).}
\label{tab:yaleb_main}
\begin{tabular}{lccc}
\hline
\textbf{Method} & \textbf{Acc.} & \textbf{Gain} \\ \hline
Forward only & 74.30\% & -- \\
Hard-switch fusion ($R_{\max}=1$) & 76.97\% & +2.67\% \\
Self-training ($R_{\max}=\infty$, 5 rounds) & \textbf{80.59\%} & \textbf{+6.29\%} \\ \hline
\end{tabular}
\end{table}

\begin{table}[t]
\centering
\caption{Results on AR (PCA-100, $K=100$, varying $n_k$, 1 trial).}
\label{tab:ar_main}
\begin{tabular}{cccc}
\hline
$n_k$ & \textbf{Fwd} & \textbf{Fus} ($R_{\max}=1$) & \textbf{ST} ($R_{\max}=2$) \\ \hline
2 & 71.33\% & 68.67\% & \textbf{82.42\%} \\
3 & 80.83\% & 79.83\% & \textbf{87.52\%} \\
4 & 88.27\% & 88.18\% & \textbf{91.68\%} \\
5 & 91.86\% & 91.86\% & \textbf{93.10\%} \\
6 & 93.80\% & 93.80\% & \textbf{95.15\%} \\ \hline
\end{tabular}
\end{table}

\subsection{Lighting Subset Evaluation}

We evaluate on Yale B's five standard lighting subsets. Training on one subset, testing on all five simultaneously (cross-lighting).

\begin{table}[t]
\centering
\caption{Single-pass fusion on lighting subsets. Train on one subset, test on all 5.}
\label{tab:lighting}
\begin{tabular}{lccccc}
\hline
\textbf{Train} & $n_k$ & \textbf{Fwd} & \textbf{Fus} & \textbf{Gain} \\ \hline
Sub1 (0--10$^\circ$) & 7 & 66.81\% & 76.73\% & +9.93\% \\
Sub2 (10--25$^\circ$) & 12 & 85.27\% & 88.77\% & +3.49\% \\
Sub3 (25--50$^\circ$) & 12 & 86.69\% & 90.64\% & +3.95\% \\
Sub4 (50--77$^\circ$) & 14 & 71.47\% & 77.47\% & +6.00\% \\
Sub5 ($>$77$^\circ$) & 19 & 85.20\% & 88.25\% & +3.04\% \\ \hline
\end{tabular}
\end{table}

All five subsets show positive gain (range +3.0\% to +9.9\%). Gains are largest when forward is weakest (Sub1, Sub4), confirming that the reverse path provides complementary evidence specifically where the inductive path fails.

\subsection{Cross-Lighting Few-Shot}

We further restrict training to only $n_k \in \{2,4,6\}$ samples from one lighting subset, testing on the remaining 62--60 images per subject across all subsets.

\begin{table}[t]
\centering
\caption{Cross-lighting few-shot. $n_k$ samples from one subset as training.}
\label{tab:fewshot}
\begin{tabular}{lcccc}
\hline
\textbf{Train} & $n_k=2$ & $n_k=4$ & $n_k=6$ \\ \hline
Sub1 & +26.3\% & +10.6\% & +8.9\% \\
Sub2 & +20.6\% & -- & +9.0\% \\
Sub3 & +16.7\% & -- & +11.3\% \\
Sub4 & +26.7\% & -- & +9.5\% \\
Sub5 & +20.7\% & -- & +10.3\% \\ \hline
\end{tabular}
\end{table}

All 25 tested configurations (5 subsets $\times$ 5 $n_k$ values) yield positive fusion gain. Gains consistently increase as $n_k$ decreases.

\subsection{Self-Training Analysis}

Self-training is most beneficial when forward accuracy is moderate (50--75\%). When forward is very weak ($<$35\%, e.g., $n_k=2$ cross-lighting), $v_j$ never exceeds 0.65 and self-training adds no samples. When forward is very strong ($>$85\%), single-pass fusion is already near-optimal and self-training risks distribution shift from pseudo-label noise.

Convergence is typically reached in 3--7 rounds. Round-by-round on Yale B ($n_k=8$):

\begin{center}
\begin{tabular}{cccc}
\hline
\textbf{Round} & \textbf{Added} & \textbf{Pseudo-label acc.} \\ \hline
1 & 1357 & 95.9\% \\
2 & 19 & 94.7\% \\
3 & 7 & 100\% \\
4 & 1 & 0\% \\
5 & 0 & converged \\ \hline
\end{tabular}
\end{center}

\subsection{Feature Space Sensitivity}

The verification gap $v_j$ is highly sensitive to feature dimensionality. In raw pixel space (2016-d), forward dominates (88.6\%) and $v_j$ gap is only 0.19, yielding negligible fusion gain (+0.04\%). In PCA-50, forward degrades to 74.3\% while reverse remains at 74.1\%, and the $v_j$ gap expands to 0.43, enabling +6.3\% gain. This pattern holds across Gabor and ResNet features, suggesting that moderate-dimensional decorrelated features are optimal for the verification mechanism.

%========================================================================
\section{Conclusion}
%========================================================================

We have shown that the reverse path in DP-GSR is fundamentally a verification mechanism, not a classifier. The verification score $v_j$ achieves a 0.43 separation gap in PCA feature spaces, enabling a simple two-threshold hard-switch fusion rule and a principled self-training pipeline. The full framework delivers consistent 3--14\% gains over forward-only baselines across two face datasets, five lighting conditions, and few-shot regimes down to 2 training samples per class, without dataset-specific parameter tuning. Future work includes integration with learned deep features and extension to open-set and domain-adaptation scenarios.

\bibliographystyle{IEEEtran}
\begin{thebibliography}{20}

\bibitem{wright2009robust}
J.~Wright, A.~Y. Yang, A.~Ganesh, S.~S. Sastry, and Y.~Ma, ``Robust face recognition via sparse representation,'' \emph{IEEE Trans. PAMI}, vol.~31, no.~2, pp. 210--227, 2009.

\bibitem{zhang2011sparse}
L.~Zhang, M.~Yang, and X.~Feng, ``Sparse representation or collaborative representation: Which helps face recognition?'' in \emph{ICCV}, 2011.

\bibitem{elhamifar2013robust}
E.~Elhamifar and R.~Vidal, ``Robust classification using structured sparse representation,'' \emph{IEEE Trans. PAMI}, vol.~35, no.~11, pp. 2773--2781, 2013.

\bibitem{cai2016probabilistic}
S.~Cai, W.~Zuo, L.~Zhang, X.~Feng, and D.~Zhang, ``A probabilistic collaborative representation based approach for fine-grained image categorization,'' in \emph{CVPR}, 2016.

\bibitem{liu2013low}
G.~Liu, Z.~Lin, S.~Yan, J.~Sun, Y.~Yu, and Y.~Ma, ``Robust recovery of subspace structures by low-rank representation,'' \emph{IEEE Trans. PAMI}, vol.~35, no.~1, pp. 171--184, 2013.

\bibitem{elhamifar2013ssc}
E.~Elhamifar and R.~Vidal, ``Sparse subspace clustering: Algorithm, theory, and applications,'' \emph{IEEE Trans. PAMI}, vol.~35, no.~11, pp. 2765--2781, 2013.

\bibitem{zhu2003semi}
X.~Zhu, Z.~Ghahramani, and J.~Lafferty, ``Semi-supervised learning using Gaussian fields and harmonic functions,'' in \emph{ICML}, 2003.

\bibitem{zhou2004learning}
D.~Zhou, O.~Bousquet, T.~N. Lal, J.~Weston, and B.~Sch{\"o}lkopf, ``Learning with local and global consistency,'' in \emph{NeurIPS}, 2004.

\bibitem{liu2018learning}
Y.~Liu, J.~Lee, M.~Park, S.~Kim, E.~Yang, S.~Hwang, and Y.~Yang, ``Learning to propagate labels: Transductive propagation network for few-shot learning,'' in \emph{ICLR}, 2019.

\bibitem{boudiaf2020information}
M.~Boudiaf, Z.~Masrourou, R.~Lidierth, and I.~Ben Ayed, ``Information maximization for few-shot learning,'' in \emph{NeurIPS}, 2020.

\bibitem{wang2024dual}
W.-Y. Wang, C.-K. Chen, X.-G. Qiao, and Y.-H. Tao, ``Dual multi-task group sparse representation for image classification,'' Preprint, 2024.

\bibitem{scudder1965probability}
H.~Scudder, ``Probability of error of some adaptive pattern-recognition machines,'' \emph{IEEE Trans. Inf. Theory}, vol.~11, no.~3, pp. 363--371, 1965.

\bibitem{georghiades2001from}
A.~S. Georghiades, P.~N. Belhumeur, and D.~J. Kriegman, ``From few to many: Illumination cone models for face recognition under variable lighting and pose,'' \emph{IEEE Trans. PAMI}, vol.~23, no.~6, pp. 643--660, 2001.

\bibitem{martinez1998ar}
A.~M. Martinez and R.~Benavente, ``The AR Face Database,'' \emph{CVC Tech. Rep.}, 1998.

\end{thebibliography}

\end{document}
